\chapter{Complexity — Universal Block}
%%% generated by the blueprint pipeline from TCSlib.Complexity.TuringMachine.UniversalBlock
%%%%%%%%%%%%%%%%%%%%%%%%%%%%%%%%%%%%%%%%%%%%%%%%%%%%%%%%%%%%%%%%%%%%%%%%%%%%%%%
\section{Overview}

This module completes the live serialized-table block of the universal Turing
machine (Arora--Barak, \S 1.4.1, Theorem~1.9): exact-cost skipping and
selection of transition records over the canonical serialized table, reading
of the eight fixed action bits, successor decoding and installation, and the
application of the decoded record, assembled into the block lemma that one
live step of the simulated machine is realized by the fixed four-tape table
interpreter within a code-dependent budget.  It also provides the concrete
checkpoint relation between source and universal-machine configurations,
its startup, halting, and output correspondence lemmas, and the block budget
itself.  Throughout, bit strings are written over $\{0,1\}$ ($1$ for true,
$0$ for false), $1^n$ denotes $n$ consecutive $1$s, and ``the table is
$w$'' means that the interpreter's first work tape is the two-way infinite
tape holding the bit string $w$ in cells $0, \dots, |w| - 1$ with blanks
elsewhere, its cursor position counted from cell $0$.

%%%%%%%%%%%%%%%%%%%%%%%%%%%%%%%%%%%%%%%%%%%%%%%%%%%%%%%%%%%%%%%%%%%%%%%%%%%%%%%
\section{Declarations}

\begin{lemma}[Exact-cost unary-tail skip]
\lean{Turing.universal_skip_unary}
\label{Turing.universal_skip_unary}
\leanok
\uses{Turing.Cfg, Turing.FinTM.bufferTape, Turing.MultiTapeTM.runFrom,
  Turing.MultiTapeTM.runFrom_succ_eq_step, Turing.MultiTapeTM.runFrom_zero,
  Turing.UniversalControl, Turing.universalEvalCfg, Turing.universalEval_step,
  Turing.universalFour, Turing.universalInterpreter, Turing.universalSkipDone}
\difficulty{4}
Suppose the table is $l \cdot 1^n \cdot 0 \cdot r$.  From the table
interpreter's unary-skipping phase, carrying a destination and a count of
remaining records in $\{0, \dots, 8\}$, with table cursor $|l|$, exactly
$n + 1$ interpreter steps advance the cursor to $|l| + n + 1$ and pass
control to the continuation after a skip request for the destination when
the remaining count is $0$, and to the fixed-field skipping phase with the
count decremented and a cleared field register otherwise; the state tape,
state cursor, and all other data are unchanged.
\end{lemma}

\begin{lemma}[Exact-cost skip of one serialized record]
\lean{Turing.universal_skip_record}
\label{Turing.universal_skip_record}
\leanok
\uses{Turing.Action, Turing.Cfg, Turing.MultiTapeTM.runFrom,
  Turing.MultiTapeTM.runFrom_add, Turing.UniversalControl,
  Turing.universalActionBits, Turing.universalEvalCfg,
  Turing.universalInterpreter, Turing.universalNextOnes,
  Turing.universalRecordBits, Turing.universalSkipDone,
  Turing.universal_record_shape, Turing.universal_skip_fixed,
  Turing.universal_skip_unary}
\difficulty{4}
Suppose the table is $l \cdot \rho \cdot r$, where $\rho$ is the serialized
transition record of a one-work-tape machine action over $\{0,1\}$ (its
eight fixed action bits followed by its self-delimiting unary successor
field).  From the interpreter's fixed-field skipping phase, carrying a
destination and a count of remaining records in $\{0, \dots, 8\}$, with
field register $0$ and table cursor $|l|$, exactly $|\rho|$ interpreter
steps advance the cursor to $|l| + |\rho|$ and pass control to the
continuation after a skip request for the destination when the remaining
count is $0$, and to the fixed-field skipping phase with the count
decremented and a cleared field register otherwise; the state tape, state
cursor, and all other data are unchanged.
\end{lemma}

\begin{lemma}[Exact-cost skip of a bounded record list]
\lean{Turing.universal_skip_records}
\label{Turing.universal_skip_records}
\leanok
\uses{Turing.Action, Turing.Cfg, Turing.MultiTapeTM.runFrom,
  Turing.MultiTapeTM.runFrom_add, Turing.UniversalControl,
  Turing.universalEvalCfg, Turing.universalInterpreter,
  Turing.universalRecordBits, Turing.universalSkipDone,
  Turing.universal_skip_record}
\difficulty{5}
Let $a_1, \dots, a_{m+1}$ be one-work-tape machine actions over $\{0,1\}$
whose serialized transition records appear consecutively in the table after
a prefix $l$, and consider the interpreter's fixed-field skipping phase
carrying a destination and the remaining-record count $m \in \{0, \dots,
8\}$, with field register $0$ and table cursor $|l|$.  Then exactly as many
interpreter steps as the total length of the $m + 1$ serialized records
advance the table cursor past all of them and enter the continuation after
a skip request for the destination; the state tape, state cursor, and all
other data are unchanged.
\end{lemma}

\begin{lemma}[Group skipping with destructive state erasure]
\lean{Turing.universal_skip_groups}
\label{Turing.universal_skip_groups}
\leanok
\uses{Turing.Action, Turing.Cfg, Turing.MultiTapeTM.runFrom,
  Turing.MultiTapeTM.runFrom_add, Turing.MultiTapeTM.runFrom_succ_eq_step,
  Turing.MultiTapeTM.runFrom_succ_eq_step', Turing.MultiTapeTM.runFrom_zero,
  Turing.UniversalControl, Turing.universalEvalCfg, Turing.universalEval_step,
  Turing.universalFour, Turing.universalInterpreter,
  Turing.universalRecordBits, Turing.universalSkipDone,
  Turing.universalStateWindow, Turing.universalStateWindow_end,
  Turing.universalStateWindow_erase, Turing.universalStateWindow_read,
  Turing.universal_skip_records}
\difficulty{5}
Suppose the table is $l \cdot G \cdot r$, where $G$ concatenates the
serialized transition records of $g$ groups of one-work-tape machine
actions, each group containing exactly nine records, and suppose the state
tape is the state window with $j$ erased and $g$ remaining unary cells,
with state cursor $j + 1$.  From the interpreter's record-group phase
carrying a read index $i \in \{0, \dots, 8\}$, with table cursor $|l|$,
exactly $g + |G| + 1$ interpreter steps reach the state-rewind phase
remembering $i$, with table cursor $|l| + |G|$, state tape the state window
with $j + g$ erased and no remaining unary cells, and state cursor
$j + g + 1$: each remaining unary state symbol is erased at the cost of one
step plus a nine-record skip, and the terminating blank costs one further
step.
\end{lemma}

\begin{lemma}[Exact-cost read of the eight fixed action bits]
\lean{Turing.universal_read_fixed}
\label{Turing.universal_read_fixed}
\leanok
\uses{Turing.Cfg, Turing.FinTM.bufferTape, Turing.FinTM.bufferTape_nat,
  Turing.MultiTapeTM.runFrom, Turing.MultiTapeTM.runFrom_succ_eq_step,
  Turing.MultiTapeTM.runFrom_zero, Turing.UniversalControl,
  Turing.universalEvalCfg, Turing.universalEval_step, Turing.universalFour,
  Turing.universalInterpreter}
\difficulty{5}
Suppose the table is $l \cdot b_0 b_1 \cdots b_7 \cdot r$ for bits
$b_0, \dots, b_7$.  Consider the interpreter's action-reading phase at a
field $f \in \{0, \dots, 7\}$, with table cursor $|l| + f$ and a bit
register agreeing with $b$ at every index below $f$.  Then exactly $8 - f$
interpreter steps reach the successor-decoding phase carrying the complete
register $b_0, \dots, b_7$, with table cursor $|l| + 8$; the state tape,
state cursor, and all other data are unchanged.
\end{lemma}

\begin{definition}[Cost of successor preparation]
\lean{Turing.universalNextCost}
\label{Turing.universalNextCost}
\leanok
The number of interpreter steps needed to decode the successor field of a
serialized record and install the successor state in unary, before the
decoded action is applied: $1$ when there is no successor state (the action
halts), and $2q + 4$ when the successor state has index $q$.
\end{definition}

\begin{lemma}[Exact-cost successor decoding and installation]
\lean{Turing.universal_prepare_next}
\label{Turing.universal_prepare_next}
\leanok
\uses{Turing.Cfg, Turing.FinTM.bufferTape, Turing.MultiTapeTM.runFrom,
  Turing.MultiTapeTM.runFrom_add, Turing.MultiTapeTM.runFrom_succ_eq_step,
  Turing.MultiTapeTM.runFrom_succ_eq_step', Turing.MultiTapeTM.runFrom_zero,
  Turing.UniversalControl, Turing.universalEvalCfg, Turing.universalEval_step,
  Turing.universalFour, Turing.universalInterpreter, Turing.universalNextCost,
  Turing.universalNextOnes, Turing.universalStateTape,
  Turing.universalStateTape_marker, Turing.universal_state_rewind,
  Turing.universal_unary_copy}
\difficulty{5}
Suppose the table is $l \cdot 1^u \cdot 0 \cdot r$, where $u$ is the unary
size of an optional successor state ($0$ when the successor is absent and
$q + 1$ when it has index $q$).  From the interpreter's successor-decoding
phase carrying eight buffered action bits, with table cursor $|l|$, state
tape the intact unary encoding of the index $0$, and state cursor $1$,
exactly the successor-preparation cost ($1$ when the successor is absent
and $2q + 4$ otherwise) in interpreter steps reaches the record-application
phase carrying the same bits and a flag recording whether the successor is
absent, with table cursor $|l| + u$, state tape the intact unary encoding
of the successor index ($0$ when absent), and state cursor $1$.
\end{lemma}

\begin{definition}[The nine actions of a coded state]
\lean{Turing.universalActions}
\label{Turing.universalActions}
\leanok
\uses{Turing.Action, Turing.CodeTM}
For a coded machine $M$ and a state $q$ of $M$, the list of the nine
transition actions of $q$: for each input read and each work read, ranging
over blank, $0$, $1$ in the fixed serialization order with the input read
varying slowest, the action $M$ performs at $q$ on that read pair.
\end{definition}

\begin{lemma}[Record lookup within a state's action list]
\lean{Turing.universalActions_lookup}
\label{Turing.universalActions_lookup}
\leanok
\uses{Turing.CodeTM, Turing.universalActions, Turing.universalHeader,
  Turing.universalRecordIndex}
\difficulty{3}
For every coded machine $M$, state $q$ of $M$, and pair of reads (an input
symbol and a work symbol, each blank, $0$, or $1$): the list of the nine
actions of $q$ has length $9$, and its entry at the record index of the
read pair (three times the serialization index of the input read plus the
serialization index of the work read) is exactly the action $M$ performs at
$q$ on those reads.
\end{lemma}

\begin{definition}[Header of the canonical serialization]
\lean{Turing.universalHeader}
\label{Turing.universalHeader}
\leanok
\uses{Turing.CodeTM, Turing.pairEncode}
The header of the canonical serialization of a coded machine $M$ with state
count $m$ and initial state index $q_0$: the pair encoding of the binary
digit string of $m$ with the empty string (each digit doubled, then the
unequal delimiter pair), followed by $1^{q_0} \cdot 0$, the initial state
in unary with its terminator.  It comprises everything preceding the
transition records.
\end{definition}

\begin{lemma}[Serialization as header and action groups]
\lean{Turing.universal_serialization_actions}
\label{Turing.universal_serialization_actions}
\leanok
\uses{Turing.CodeTM, Turing.CodeTM.serialize, Turing.pairEncode,
  Turing.universalActions, Turing.universalHeader,
  Turing.universalRecordBits, Turing.universalRecords}
\difficulty{5}
The canonical serialization of a coded machine $M$ equals its header
followed by the concatenation, over every state of $M$ in increasing order,
of the concatenated serialized transition records of that state's
nine-action list (its actions on all read pairs, in the fixed input-major
serialization order).
\end{lemma}

\begin{lemma}[Decomposition of the table at a selected record]
\lean{Turing.universal_lookup_parts}
\label{Turing.universal_lookup_parts}
\leanok
\uses{Turing.Action, Turing.CodeTM, Turing.CodeTM.serialize,
  Turing.universalActions, Turing.universalActions_lookup,
  Turing.universalHeader, Turing.universalRecordBits,
  Turing.universalRecordIndex, Turing.universal_serialization_actions}
\difficulty{5}
For every coded machine $M$, state $q$ of $M$, and pair of reads (an input
symbol and a work symbol, each blank, $0$, or $1$), the canonical
serialization of $M$ decomposes as: the header, then the concatenated
serialized records of a list of exactly $q$ groups of nine actions each
(the full groups of the states preceding $q$), then the concatenated
serialized records of a list of actions whose length is the record index of
the read pair, then the serialized record of the action $M$ performs at $q$
on those reads, then a remainder.
\end{lemma}

\begin{lemma}[Decoding the eight fixed bits]
\lean{Turing.universalActionBits_decode}
\label{Turing.universalActionBits_decode}
\leanok
\uses{Turing.Action, Turing.universalActionBits, Turing.universalSign,
  Turing.universalWrite}
\difficulty{3}
For every one-work-tape machine action $a$ over $\{0,1\}$, the eight fixed
bits of its serialized record decode back to its components: bits $0, 1$
decode to $a$'s input-head movement, bits $2, 3$ to its optional work-tape
write, bits $4, 5$ to its work-head movement, and bits $6, 7$ to its
optional output emission (bit $6$ flags whether an output is emitted and
bit $7$ is the emitted bit).
\end{lemma}

\begin{lemma}[One step applies the decoded record]
\lean{Turing.universal_apply_record}
\label{Turing.universal_apply_record}
\leanok
\uses{Turing.Action, Turing.Action.apply, Turing.Cfg, Turing.Cfg.inputSymbol,
  Turing.Cfg.workTapeSymbols, Turing.CodeTM, Turing.CodeTM.serialize,
  Turing.FinTM.bufferTape, Turing.FinTM.virtualMove, Turing.MultiTapeTM.step,
  Turing.UniversalControl, Turing.universalActionBits,
  Turing.universalActionBits_decode, Turing.universalEvalCfg,
  Turing.universalFour, Turing.universalInput_move,
  Turing.universalInput_read, Turing.universalInterpreter,
  Turing.universalSimulationCfg, Turing.universalStateTape}
\difficulty{5}
Let $M$ be a coded machine, $\alpha$ a code string, $s$ a configuration of
$M$'s one-work-tape form on an input $x$, and $a$ a transition action of
$M$.  One step of the table interpreter from the evaluation configuration
over the interpreter checkpoint of $s$ (at an arbitrary table cursor) whose
control is the record-application phase carrying the eight fixed bits of
$a$ and the flag recording whether $a$ halts, whose table is the canonical
serialization of $M$ at a cursor $p$, and whose state tape is the intact
unary encoding of $a$'s successor index ($0$ when $a$ halts) with state
cursor $1$, is exactly the interpreter checkpoint of the configuration
obtained by applying $a$ to $s$, with table cursor $p$.
\end{lemma}

\begin{lemma}[Exact-cost record selection]
\lean{Turing.universal_select}
\label{Turing.universal_select}
\leanok
\uses{Turing.Action, Turing.Cfg, Turing.MultiTapeTM.runFrom,
  Turing.MultiTapeTM.runFrom_add, Turing.UniversalControl,
  Turing.universalEvalCfg, Turing.universalInterpreter,
  Turing.universalRecordBits, Turing.universalSkipDone,
  Turing.universalStateTape, Turing.universalStateTape_marker,
  Turing.universalStateWindow_empty, Turing.universalStateWindow_zero,
  Turing.universal_skip_groups, Turing.universal_skip_records,
  Turing.universal_state_rewind}
\difficulty{5}
Suppose the table is $l \cdot G \cdot B \cdot r$, where $G$ concatenates
the serialized transition records of $g$ groups of nine one-work-tape
machine actions each and $B$ concatenates the serialized records of a list
of actions whose length is a read index $i \in \{0, \dots, 8\}$, and
suppose the state tape is the intact unary encoding of $g$ with state
cursor $1$.  From the interpreter's record-group phase carrying $i$, with
table cursor $|l|$, exactly $2g + |G| + |B| + 3$ interpreter steps reach
the action-reading phase at field $0$ with a cleared bit register, table
cursor $|l| + |G| + |B|$, state tape the intact unary encoding of $0$, and
state cursor $1$: the two table scans cost their serialized lengths, and
the destructive state count with its rewind costs $2g + 3$.
\end{lemma}

\begin{lemma}[Concatenation of exact runs]
\lean{Turing.universal_run_join}
\label{Turing.universal_run_join}
\leanok
\uses{Turing.Cfg, Turing.MultiTapeTM, Turing.MultiTapeTM.runFrom,
  Turing.MultiTapeTM.runFrom_add}
\difficulty{1}
For a multi-tape machine over $\{0,1\}$ and configurations $a$, $b$, $c$ on
a common input: if running the machine for $s$ steps from $a$ yields $b$,
and running it for $t$ steps from $b$ yields $c$, then running it for
$s + t$ steps from $a$ yields $c$.
\end{lemma}

\begin{lemma}[One live source step within the block bound]
\lean{Turing.universal_live_block}
\label{Turing.universal_live_block}
\leanok
\uses{Turing.Action.apply, Turing.Cfg, Turing.Cfg.inputSymbol,
  Turing.Cfg.workTapeSymbols, Turing.CodeTM, Turing.CodeTM.serialize,
  Turing.MultiTapeTM.runFrom, Turing.MultiTapeTM.runFrom_succ_eq_step',
  Turing.MultiTapeTM.runFrom_zero, Turing.MultiTapeTM.step,
  Turing.pairEncode, Turing.universalActionBits, Turing.universalAdmin,
  Turing.universalBlockBound, Turing.universalEvalCfg,
  Turing.universalEval_step, Turing.universalHeader,
  Turing.universalInput_read, Turing.universalInterpreter,
  Turing.universalNextCost, Turing.universalNextOnes,
  Turing.universalRecordBits, Turing.universalRecordIndex,
  Turing.universalSimulationCfg, Turing.universalStateTape,
  Turing.universal_apply_record, Turing.universal_count_run,
  Turing.universal_initial_skip, Turing.universal_lookup_parts,
  Turing.universal_pair_length, Turing.universal_prepare_next,
  Turing.universal_read_fixed, Turing.universal_record_shape,
  Turing.universal_run_join, Turing.universal_select,
  Turing.universal_table_rewind}
\difficulty{6}
Let $M$ be a coded machine with state count $N$ and canonical serialization
of length $L$, let $\alpha$ be a code string, let $s$ be a configuration of
$M$'s one-work-tape form on an input $x$ that has not halted, and let
$p \le L$ be a table cursor.  Then there exist a duration $d$ with
$1 \le d \le 3L + 5(N + 1) + 20$ and a cursor $p' \le L$ such that running
the table interpreter for $d$ steps from the interpreter checkpoint of $s$
with table cursor $p$ yields exactly the interpreter checkpoint of the
configuration obtained from $s$ by one step of $M$, with table cursor
$p'$.
\end{lemma}

\begin{definition}[The block budget]
\lean{Turing.universalBlockBound}
\label{Turing.universalBlockBound}
\leanok
\uses{Turing.CodeTM.serialize, Turing.EffectiveMachineCode}
The code-dependent block budget of the universal machine for an effective
machine-coding scheme on a code string $\alpha$:
\[
  3L + 5(N + 1) + 20,
\]
where $L$ is the length of the canonical serialization of the machine
decoded from $\alpha$ and $N$ is that machine's state count.
\end{definition}

\begin{definition}[The concrete checkpoint relation]
\lean{Turing.universalRelation}
\label{Turing.universalRelation}
\leanok
\uses{Turing.Cfg, Turing.CodeTM.serialize, Turing.EffectiveMachineCode,
  Turing.FinTM.rightCfg, Turing.pairEncode, Turing.universalCanonTM,
  Turing.universalSimulationCfg, Turing.universalTM}
For an effective machine-coding scheme $c$, a code string $\alpha$, and an
input $x$: the relation between a configuration $s$ of the machine decoded
from $\alpha$ (in one-work-tape form) on $x$ and a configuration $t$ of the
complete universal-machine candidate on the pair encoding
$\langle \alpha, x \rangle$ that holds when there exist a table cursor $p$
not exceeding the length of the canonical serialization of the decoded
machine, together with contents and head positions for the inactive
canonizer tapes, such that $t$ is the right-block embedding, over those
inactive tapes, of the interpreter checkpoint of $s$ with table cursor
$p$.
\end{definition}

\begin{lemma}[The header ends within the table]
\lean{Turing.universal_header_bound}
\label{Turing.universal_header_bound}
\leanok
\uses{Turing.CodeTM, Turing.CodeTM.serialize, Turing.universal_pair_length,
  Turing.universal_serialization_header}
\difficulty{3}
For every coded machine $M$ with state count $m$ and initial state index
$q_0$: writing $\ell$ for the number of binary digits of $m$ and $L$ for
the length of the canonical serialization of $M$, one has
$2\ell + 2 + q_0 + 1 \le L$; that is, the position just past the doubled
count field, its delimiter pair, and the unary initial-state field with its
terminator lies within the serialized table.
\end{lemma}

\begin{lemma}[Startup establishes the checkpoint relation]
\lean{Turing.universalRelation_start}
\label{Turing.universalRelation_start}
\leanok
\uses{Turing.EffectiveMachineCode, Turing.MultiTapeTM.initCfg,
  Turing.MultiTapeTM.runFrom, Turing.pairEncode, Turing.universalRelation,
  Turing.universalStartupBound, Turing.universalTM,
  Turing.universal_header_bound, Turing.universal_initialized}
\difficulty{1}
For every effective machine-coding scheme $c$, code string $\alpha$, and
input $x$, there is a time $t$ not exceeding the startup bound for $\alpha$
such that the concrete checkpoint relation holds between the initial
configuration of the machine decoded from $\alpha$ on $x$ and the
configuration of the complete universal-machine candidate after $t$ steps
from its initial configuration on the pair encoding
$\langle \alpha, x \rangle$.
\end{lemma}

\begin{lemma}[The checkpoint relation preserves halting]
\lean{Turing.universalRelation_halt}
\label{Turing.universalRelation_halt}
\leanok
\uses{Turing.Cfg, Turing.EffectiveMachineCode, Turing.FinTM.rightCfg,
  Turing.pairEncode, Turing.universalRelation,
  Turing.universalSimulationCfg, Turing.universalTM}
\difficulty{1}
Let $c$ be an effective machine-coding scheme, $\alpha$ a code string, and
$x$ an input.  If the concrete checkpoint relation holds between a
configuration $s$ of the machine decoded from $\alpha$ on $x$ and a
configuration $t$ of the complete universal-machine candidate on the pair
encoding $\langle \alpha, x \rangle$, then $s$ has halted if and only if
$t$ has halted.
\end{lemma}

\begin{lemma}[The checkpoint relation preserves output]
\lean{Turing.universalRelation_output}
\label{Turing.universalRelation_output}
\leanok
\uses{Turing.Cfg, Turing.EffectiveMachineCode, Turing.pairEncode,
  Turing.universalRelation, Turing.universalTM}
\difficulty{1}
Let $c$ be an effective machine-coding scheme, $\alpha$ a code string, and
$x$ an input.  If the concrete checkpoint relation holds between a
configuration $s$ of the machine decoded from $\alpha$ on $x$ and a
configuration $t$ of the complete universal-machine candidate on the pair
encoding $\langle \alpha, x \rangle$, then the accumulated output of $s$
equals the accumulated output of $t$.
\end{lemma}
